\documentclass[11pt,a4paper]{article}
\usepackage[T1]{fontenc}
\usepackage[utf8]{inputenc}
\usepackage[english]{babel}
\usepackage{amsmath,amssymb,mathtools}
\usepackage[margin=26mm]{geometry}
\usepackage[colorlinks=true,linkcolor=blue,urlcolor=blue,citecolor=blue]{hyperref}
\title{From Mathematical Exploration to a Candidate Proof:\\
An Experiment in AI-Assisted Research on the\\
Three-Dimensional Covariogram Problem}
\author{Author details to be supplied}
\date{9 October 2026 --- proposed structure for author approval}

\begin{document}
\maketitle

\noindent\textbf{Editorial status.}
This document proposes the title, abstract, section structure, and accompanying
materials. It is not the completed paper. The full manuscript and public
repository will be prepared after the author has approved this structure.
The existing mathematical notebook remains a separate supporting document.

\begin{abstract}
We document an experiment in AI-assisted mathematical research led by an
author who is neither a specialist in convex geometry nor a professional
researcher in the field. Starting from the three-dimensional covariogram
problem, the investigation combined literature study, exploratory arguments,
symbolic and numerical diagnostics, and partial formalization in Lean. Its
principal mathematical artifact is a candidate argument for the determination
of a $C^9_+$ convex body in $\mathbb R^3$ among arbitrary convex competitors,
up to translation and central reflection. The argument remains subject to
expert review. Thirty-three supporting lemmas and conditional deductions have
been checked in Lean; the complete uniqueness statement and essential
geometric and spectral links have not been formalized. The contribution of
this report is a documented research process and an inspectable collection
of mathematical artifacts, rather than an established solution of the open
case. We examine the opportunities and limits illustrated by this experience
and ask what domain experts might achieve by combining these tools with
specialist judgment. A single exploratory case cannot establish general
productivity gains, correctness, or mathematical novelty, but it motivates
systematic investigation of this research methodology.
\end{abstract}

\section{Introduction: Research Assistance Beyond Specialist Expertise}
\paragraph{Purpose.}
Explain the author's position, the motivation for attempting a difficult
problem outside the author's professional research field, and the methodological
question that motivates the paper. Describe the experiment as a case study,
without implying that specialist knowledge has become unnecessary.

\paragraph{Central question.}
\emph{If an author without specialist expertise can use AI tools to develop
an inspectable candidate argument and formally check some of its components,
what might a domain expert achieve by combining the same tools with deeper
knowledge and stronger mathematical judgment?}

The introduction will distinguish observed outcomes from expectations about
future tools and expert practice. Related work on AI-assisted mathematics and
formal proof will be added after a focused literature review; no claim of
priority for the methodology is made in this outline.

\section{The Mathematical Case: A Precise Covariogram Uniqueness Question}
Introduce a convex body as a compact convex subset of $\mathbb R^3$ with
nonempty interior, and define
\[
 g_K(x)=\operatorname{Vol}_3\bigl(K\cap(K+x)\bigr).
\]
Explain translation and central-reflection invariance. State the target
precisely:
\[
\begin{gathered}
 K\in C^9_+,\qquad L\subset\mathbb R^3\text{ an arbitrary convex body},
 \qquad g_K=g_L,\\
 \Longrightarrow\quad
 \exists v\in\mathbb R^3:\quad L=v+K\ \text{or}\ L=v-K.
\end{gathered}
\]
Here $C^9_+$ denotes $C^9$ boundary regularity with positive Gaussian curvature;
no boundary regularity is imposed on $L$. Locate this question relative to
Bianchi's results and the extension discussed in his survey
\cite{Bianchi2016,Bianchi2023}. Distinguish the historical statement that the
case was reported open from any claim about its current status or the novelty
of the present candidate.

\section{Experimental Setting and Human--AI Collaboration}
Describe the actual tools and available records, the human choice of problems
and constraints, and the AI's contributions to literature exploration,
argument generation, criticism, computation, formalization, and writing.
Separate recorded actions from retrospective interpretation.

Report model and software versions only where recoverable, with any missing
metadata identified explicitly. Explain what can be reproduced from the
repository and what cannot be recreated from an interactive session. The
experiment is observational, not a controlled comparison of human and AI
performance.

\section{Development of the Candidate Argument}
Reconstruct the investigation from the broad three-dimensional problem to
the target with one regular body and an arbitrary competitor. Include the
main explored approaches, discarded or insufficient arguments, changes in
scope, and reasons for selecting the final route.

Identify concrete benefits of the workflow that are supported by the records:
organizing dependencies, producing explicit calculations, constructing
diagnostics, and exposing the difference between an informal argument and
formal verification. Do not convert the apparent fluency or length of an
argument into evidence of correctness.

\section{The Candidate Proof and Its Dependencies}
Give a readable account of the mathematical argument, with the complete
proposed derivations in Appendix~\ref{app:proof}. Use explicit status labels
for cited results, proposed lemmas, checked components, and remaining
verification obligations.

\subsection{Information Recovered from the Covariogram}
Describe the proposed chain from equal covariograms to the difference body,
projection data, and identities for Fourier--Laplace transforms. State the
external results used and their exact hypotheses.

\subsection{The Arbitrary Competitor: Support Regularity and Area Densities}
Present the proposed lemma concerning the support function of $L$, its
almost-everywhere quadratic data, the absence of a singular part of the area
measure, and endpoint section asymptotics, including degenerate cases.
Flag this as a principal target for expert scrutiny.

\subsection{Fourier--Laplace Zeros and a Common Orientation Parameter}
Describe the use of Bianchi's zero branches, factor allocation, Hadamard
factorization, logarithmic kernels, and weighted limits. Separate the
formalized analytic estimates from the unformalized derivation of their
hypotheses for the geometric problem.

\subsection{Conditional Completion Through Area Measures}
Explain the proposed recovery of globally ordered area densities and the
conclusion via Minkowski uniqueness. Do not describe this final chain as
established merely because some of its conditional steps have been checked.

\section{Verification: Evidence, Formal Coverage, and Remaining Obligations}
\subsection{Symbolic and Numerical Diagnostics}
Describe the rational-arithmetic checks and the truncated numerical
experiments, including precisely what they test and what they do not certify.
Preserve the distinction between numerical agreement and proof.

\subsection{Partial Verification in Lean}
Document the 33 explicit theorems in nine modules, their hypotheses, the
Lean/mathlib versions, and the successful compilation and axiom audit.
Explain that the usual foundational axioms are used, with no admitted proofs
or additional mathematical axioms. Include the verified logarithmic
remainder estimates, convergence arguments, and conclusions about densities.

\subsection{What Has Not Been Verified}
List the missing geometric and spectral links, the full factorization and
Laplace argument, the application to the specific weights, and the geometric
uniqueness conclusion. State prominently that neither central candidate
lemma is certified in its entirety and that theorem counts do not measure
distance to a complete proof.

\section{Methodological Lessons and Prospects for Expert Research}
Discuss the possible role of AI tools in lowering the cost of exploration,
organizing proof obligations, and connecting informal arguments to formal
proof assistants. Interpret these as possibilities illustrated by the case,
not measured general improvements.

Explain where expert knowledge may be especially valuable: detecting missing
hypotheses, recognizing existing results, selecting feasible approaches, and
rejecting persuasive but invalid arguments. Address limitations, including
selection bias, the absence of a comparison group, incomplete interaction
records, and the time required for independent validation.

\section{Reproducibility, Expert Review, and Next Steps}
Specify the public repository contents, the commands for repeating the
checks, and the correspondence between paper claims and source files.
The present Lean setup uses local tool and library paths; portable setup
instructions are a release requirement, not an already completed feature.

Set out the intended review process: independent assessment of the candidate
lemmas, further verification of bibliographic status, and extension of the
formalization. Record any future corrections and expert feedback transparently.
A planned request for Bianchi's assessment must not be represented as an
endorsement or as a review already received.

\section{Conclusion}
Summarize the methodological artifact produced by the experiment and the
precise mathematical status of the candidate. Return to the question of what
specialists might accomplish with these tools. The methodological record
remains informative if expert review later requires substantial corrections
or rejects the candidate; it must not depend on advertising an unverified
solution.

\section*{Author Contributions and Disclosure of AI Assistance}
Describe the human author's role and the substantial assistance used in
research, drafting, computation, and formalization. Identify tools and models
where the records support doing so. State human responsibility for the
manuscript and its claims. Author name, affiliation or independent status,
repository URL, acknowledgments, and licenses remain to be supplied or agreed
before release.

\appendix
\section{Detailed Candidate Argument}\label{app:proof}
Provide the English technical account with all proposed derivations,
dependencies, and unresolved verification obligations. Preserve the original
Italian notebook as a separate, dated research artifact rather than silently
replacing it with a polished retrospective account.

\section{Lean Theorems and Coverage Map}
Map each theorem name and module to the corresponding mathematical step,
including its premises. Include the audit output, version information, and
source checksums. Distinguish rerunning the verification from establishing
that the formal statements cover the intended uniqueness theorem.

\section{Computational Diagnostics}
Document the symbolic and numerical scripts, inputs, outputs, truncation
choices, and limitations. Relate them to the exact statements investigated
in the notebook and the paper.

\section{Repository Guide and Research Record}
The release is intended to contain the following materials:
\begin{itemize}
\item The English paper source and a compiled reading copy.
\item The original mathematical notebook and its verification notes.
\item The Lean sources, pinned dependencies, portable setup instructions,
      build script, axiom audit, and verification logs.
\item The rational and numerical diagnostic scripts with instructions.
\item The bibliography register, a record of the available research stages,
      and a correction history.
\end{itemize}
Downloaded toolchains, dependency caches, machine-specific absolute paths,
and private session metadata are not intended as public repository content.
Any unavailable interaction history will be reported as unavailable, not
reconstructed as though it were a verbatim record.

\begin{thebibliography}{9}
\bibitem{Bianchi2016}
G.~Bianchi, \emph{The covariogram and Fourier--Laplace transform in
$\mathbb C^n$}, Proceedings of the London Mathematical Society
113 (2016), 1--23.
\href{https://arxiv.org/abs/1312.7816}{arXiv:1312.7816}.

\bibitem{Bianchi2023}
G.~Bianchi, \emph{The covariogram problem}, in
\emph{Harmonic Analysis and Convexity} (2023), 37--82.
\href{https://people.dimai.unifi.it/bianchi/lavori/covariogram_survey14.pdf}{Author's manuscript}.
\end{thebibliography}
\end{document}
